% Options for packages loaded elsewhere
\PassOptionsToPackage{unicode}{hyperref}
\PassOptionsToPackage{hyphens}{url}
\documentclass[
]{article}
\usepackage{xcolor}
\usepackage[margin=1in]{geometry}
\usepackage{amsmath,amssymb}
\setcounter{secnumdepth}{-\maxdimen} % remove section numbering
\usepackage{iftex}
\ifPDFTeX
  \usepackage[T1]{fontenc}
  \usepackage[utf8]{inputenc}
  \usepackage{textcomp} % provide euro and other symbols
\else % if luatex or xetex
  \usepackage{unicode-math} % this also loads fontspec
  \defaultfontfeatures{Scale=MatchLowercase}
  \defaultfontfeatures[\rmfamily]{Ligatures=TeX,Scale=1}
\fi
\usepackage{lmodern}
\ifPDFTeX\else
  % xetex/luatex font selection
\fi
% Use upquote if available, for straight quotes in verbatim environments
\IfFileExists{upquote.sty}{\usepackage{upquote}}{}
\IfFileExists{microtype.sty}{% use microtype if available
  \usepackage[]{microtype}
  \UseMicrotypeSet[protrusion]{basicmath} % disable protrusion for tt fonts
}{}
\makeatletter
\@ifundefined{KOMAClassName}{% if non-KOMA class
  \IfFileExists{parskip.sty}{%
    \usepackage{parskip}
  }{% else
    \setlength{\parindent}{0pt}
    \setlength{\parskip}{6pt plus 2pt minus 1pt}}
}{% if KOMA class
  \KOMAoptions{parskip=half}}
\makeatother
\usepackage{color}
\usepackage{fancyvrb}
\newcommand{\VerbBar}{|}
\newcommand{\VERB}{\Verb[commandchars=\\\{\}]}
\DefineVerbatimEnvironment{Highlighting}{Verbatim}{commandchars=\\\{\}}
% Add ',fontsize=\small' for more characters per line
\usepackage{framed}
\definecolor{shadecolor}{RGB}{248,248,248}
\newenvironment{Shaded}{\begin{snugshade}}{\end{snugshade}}
\newcommand{\AlertTok}[1]{\textcolor[rgb]{0.94,0.16,0.16}{#1}}
\newcommand{\AnnotationTok}[1]{\textcolor[rgb]{0.56,0.35,0.01}{\textbf{\textit{#1}}}}
\newcommand{\AttributeTok}[1]{\textcolor[rgb]{0.13,0.29,0.53}{#1}}
\newcommand{\BaseNTok}[1]{\textcolor[rgb]{0.00,0.00,0.81}{#1}}
\newcommand{\BuiltInTok}[1]{#1}
\newcommand{\CharTok}[1]{\textcolor[rgb]{0.31,0.60,0.02}{#1}}
\newcommand{\CommentTok}[1]{\textcolor[rgb]{0.56,0.35,0.01}{\textit{#1}}}
\newcommand{\CommentVarTok}[1]{\textcolor[rgb]{0.56,0.35,0.01}{\textbf{\textit{#1}}}}
\newcommand{\ConstantTok}[1]{\textcolor[rgb]{0.56,0.35,0.01}{#1}}
\newcommand{\ControlFlowTok}[1]{\textcolor[rgb]{0.13,0.29,0.53}{\textbf{#1}}}
\newcommand{\DataTypeTok}[1]{\textcolor[rgb]{0.13,0.29,0.53}{#1}}
\newcommand{\DecValTok}[1]{\textcolor[rgb]{0.00,0.00,0.81}{#1}}
\newcommand{\DocumentationTok}[1]{\textcolor[rgb]{0.56,0.35,0.01}{\textbf{\textit{#1}}}}
\newcommand{\ErrorTok}[1]{\textcolor[rgb]{0.64,0.00,0.00}{\textbf{#1}}}
\newcommand{\ExtensionTok}[1]{#1}
\newcommand{\FloatTok}[1]{\textcolor[rgb]{0.00,0.00,0.81}{#1}}
\newcommand{\FunctionTok}[1]{\textcolor[rgb]{0.13,0.29,0.53}{\textbf{#1}}}
\newcommand{\ImportTok}[1]{#1}
\newcommand{\InformationTok}[1]{\textcolor[rgb]{0.56,0.35,0.01}{\textbf{\textit{#1}}}}
\newcommand{\KeywordTok}[1]{\textcolor[rgb]{0.13,0.29,0.53}{\textbf{#1}}}
\newcommand{\NormalTok}[1]{#1}
\newcommand{\OperatorTok}[1]{\textcolor[rgb]{0.81,0.36,0.00}{\textbf{#1}}}
\newcommand{\OtherTok}[1]{\textcolor[rgb]{0.56,0.35,0.01}{#1}}
\newcommand{\PreprocessorTok}[1]{\textcolor[rgb]{0.56,0.35,0.01}{\textit{#1}}}
\newcommand{\RegionMarkerTok}[1]{#1}
\newcommand{\SpecialCharTok}[1]{\textcolor[rgb]{0.81,0.36,0.00}{\textbf{#1}}}
\newcommand{\SpecialStringTok}[1]{\textcolor[rgb]{0.31,0.60,0.02}{#1}}
\newcommand{\StringTok}[1]{\textcolor[rgb]{0.31,0.60,0.02}{#1}}
\newcommand{\VariableTok}[1]{\textcolor[rgb]{0.00,0.00,0.00}{#1}}
\newcommand{\VerbatimStringTok}[1]{\textcolor[rgb]{0.31,0.60,0.02}{#1}}
\newcommand{\WarningTok}[1]{\textcolor[rgb]{0.56,0.35,0.01}{\textbf{\textit{#1}}}}
\usepackage{graphicx}
\makeatletter
\newsavebox\pandoc@box
\newcommand*\pandocbounded[1]{% scales image to fit in text height/width
  \sbox\pandoc@box{#1}%
  \Gscale@div\@tempa{\textheight}{\dimexpr\ht\pandoc@box+\dp\pandoc@box\relax}%
  \Gscale@div\@tempb{\linewidth}{\wd\pandoc@box}%
  \ifdim\@tempb\p@<\@tempa\p@\let\@tempa\@tempb\fi% select the smaller of both
  \ifdim\@tempa\p@<\p@\scalebox{\@tempa}{\usebox\pandoc@box}%
  \else\usebox{\pandoc@box}%
  \fi%
}
% Set default figure placement to htbp
\def\fps@figure{htbp}
\makeatother
\setlength{\emergencystretch}{3em} % prevent overfull lines
\providecommand{\tightlist}{%
  \setlength{\itemsep}{0pt}\setlength{\parskip}{0pt}}
\usepackage{bookmark}
\IfFileExists{xurl.sty}{\usepackage{xurl}}{} % add URL line breaks if available
\urlstyle{same}
\hypersetup{
  pdftitle={Probability},
  pdfauthor={Guibril Ramde},
  hidelinks,
  pdfcreator={LaTeX via pandoc}}

\title{Probability}
\author{Guibril Ramde}
\date{}

\begin{document}
\maketitle

\subsection{The Hot Hand}\label{the-hot-hand}

Basketball players who make several baskets in succession are described
as having a \emph{hot hand}. Fans and players have long believed in the
hot hand phenomenon, which refutes the assumption that each shot is
independent of the next. However,
\href{http://www.sciencedirect.com/science/article/pii/0010028585900106}{a
1985 paper} by Gilovich, Vallone, and Tversky collected evidence that
contradicted this belief and showed that successive shots are
independent events. This paper started a great controversy that
continues to this day, as you can see by Googling \emph{hot hand
basketball}.

We do not expect to resolve this controversy today. However, in this lab
we'll apply one approach to answering questions like this. The goals for
this lab are to (1) think about the effects of independent and dependent
events, (2) learn how to simulate shooting streaks in R, and (3) to
compare a simulation to actual data in order to determine if the hot
hand phenomenon appears to be real.

\subsection{Getting Started}\label{getting-started}

\subsubsection{Load packages}\label{load-packages}

In this lab, we will explore and visualize the data using the
\texttt{tidyverse} suite of packages. The data can be found in the
companion package for OpenIntro labs, \textbf{openintro}.

Let's load the packages.

\begin{Shaded}
\begin{Highlighting}[]
\FunctionTok{library}\NormalTok{(tidyverse)}
\FunctionTok{library}\NormalTok{(openintro)}
\end{Highlighting}
\end{Shaded}

\subsubsection{Data}\label{data}

Your investigation will focus on the performance of one player:
\href{https://en.wikipedia.org/wiki/Kobe_Bryant}{Kobe Bryant} of the Los
Angeles Lakers. His performance against the Orlando Magic in the
\href{https://en.wikipedia.org/wiki/2009_NBA_Finals}{2009 NBA Finals}
earned him the title \emph{Most Valuable Player} and many spectators
commented on how he appeared to show a hot hand. The data file we'll use
is called \texttt{kobe\_basket}.

\begin{Shaded}
\begin{Highlighting}[]
\FunctionTok{glimpse}\NormalTok{(kobe\_basket)}
\end{Highlighting}
\end{Shaded}

\begin{verbatim}
## Rows: 133
## Columns: 6
## $ vs          <fct> ORL, ORL, ORL, ORL, ORL, ORL, ORL, ORL, ORL, ORL, ORL, ORL~
## $ game        <int> 1, 1, 1, 1, 1, 1, 1, 1, 1, 1, 1, 1, 1, 1, 1, 1, 1, 1, 1, 1~
## $ quarter     <fct> 1, 1, 1, 1, 1, 1, 1, 1, 1, 2, 2, 2, 2, 2, 2, 2, 2, 2, 3, 3~
## $ time        <fct> 9:47, 9:07, 8:11, 7:41, 7:03, 6:01, 4:07, 0:52, 0:00, 6:35~
## $ description <fct> Kobe Bryant makes 4-foot two point shot, Kobe Bryant misse~
## $ shot        <chr> "H", "M", "M", "H", "H", "M", "M", "M", "M", "H", "H", "H"~
\end{verbatim}

This data frame contains 133 observations and 6 variables, where every
row records a shot taken by Kobe Bryant. The \texttt{shot} variable in
this dataset indicates whether the shot was a hit (\texttt{H}) or a miss
(\texttt{M}).

Just looking at the string of hits and misses, it can be difficult to
gauge whether or not it seems like Kobe was shooting with a hot hand.
One way we can approach this is by considering the belief that hot hand
shooters tend to go on shooting streaks. For this lab, we define the
length of a shooting streak to be the \emph{number of consecutive
baskets made until a miss occurs}.

For example, in Game 1 Kobe had the following sequence of hits and
misses from his nine shot attempts in the first quarter:

\[ \textrm{H M | M | H H M | M | M | M} \]

You can verify this by viewing the first 9 rows of the data in the data
viewer.

Within the nine shot attempts, there are six streaks, which are
separated by a ``\textbar{}'' above. Their lengths are one, zero, two,
zero, zero, zero (in order of occurrence).

\begin{enumerate}
\def\labelenumi{\arabic{enumi}.}
\tightlist
\item
  What does a streak length of 1 mean, i.e.~how many hits and misses are
  in a streak of 1? What about a streak length of 0?
\end{enumerate}

\textbf{Insert your answer here} A streak length of 1 means Kobe made
one shot and then missed the next shot. Therefore, a streak of 1
consists of one hit followed by one miss.

A streak length of 0 means Kobe missed immediately without making any
shots first. In this case, there are zero hits before the miss.

Counting streak lengths manually for all 133 shots would get tedious, so
we'll use the custom function \texttt{calc\_streak} to calculate them,
and store the results in a data frame called \texttt{kobe\_streak} as
the \texttt{length} variable.

\begin{Shaded}
\begin{Highlighting}[]
\NormalTok{kobe\_streak }\OtherTok{\textless{}{-}}\NormalTok{ openintro}\SpecialCharTok{::}\FunctionTok{calc\_streak}\NormalTok{(kobe\_basket}\SpecialCharTok{$}\NormalTok{shot)}
\NormalTok{kobe\_streak}
\end{Highlighting}
\end{Shaded}

\begin{verbatim}
##    length
## 1       1
## 2       0
## 3       2
## 4       0
## 5       0
## 6       0
## 7       3
## 8       2
## 9       0
## 10      3
## 11      0
## 12      1
## 13      3
## 14      0
## 15      0
## 16      0
## 17      0
## 18      0
## 19      1
## 20      1
## 21      0
## 22      4
## 23      1
## 24      0
## 25      1
## 26      0
## 27      1
## 28      0
## 29      1
## 30      2
## 31      0
## 32      1
## 33      2
## 34      1
## 35      0
## 36      0
## 37      1
## 38      0
## 39      0
## 40      0
## 41      1
## 42      1
## 43      0
## 44      1
## 45      0
## 46      2
## 47      0
## 48      0
## 49      0
## 50      3
## 51      0
## 52      1
## 53      0
## 54      1
## 55      2
## 56      1
## 57      0
## 58      1
## 59      0
## 60      0
## 61      1
## 62      3
## 63      3
## 64      1
## 65      1
## 66      0
## 67      0
## 68      0
## 69      0
## 70      0
## 71      1
## 72      1
## 73      0
## 74      0
## 75      0
## 76      1
\end{verbatim}

We can then take a look at the distribution of these streak lengths.

\begin{Shaded}
\begin{Highlighting}[]
\FunctionTok{ggplot}\NormalTok{(}\AttributeTok{data =}\NormalTok{ kobe\_streak, }\FunctionTok{aes}\NormalTok{(}\AttributeTok{x =}\NormalTok{ length)) }\SpecialCharTok{+}
  \FunctionTok{geom\_bar}\NormalTok{()}
\end{Highlighting}
\end{Shaded}

\pandocbounded{\includegraphics[keepaspectratio]{Lab3_probability_files/figure-latex/plot-streak-kobe-1.pdf}}

\begin{enumerate}
\def\labelenumi{\arabic{enumi}.}
\setcounter{enumi}{1}
\tightlist
\item
  Describe the distribution of Kobe's streak lengths from the 2009 NBA
  finals. What was his typical streak length? How long was his longest
  streak of baskets? Make sure to include the accompanying plot in your
  answer.
\end{enumerate}

\textbf{Insert your answer here} The distribution of Kobe's streak
lengths is right-skewed. Most of the streaks are short (0 or 1), and the
frequency decreases as the streak length increases. This means short
streaks occurred much more often than long streaks. his typical streak
length is 1. the longest streak of baskets is 4.

\begin{Shaded}
\begin{Highlighting}[]
\FunctionTok{ggplot}\NormalTok{(}\AttributeTok{data =}\NormalTok{ kobe\_streak, }\FunctionTok{aes}\NormalTok{(}\AttributeTok{x =}\NormalTok{ length)) }\SpecialCharTok{+} \FunctionTok{geom\_bar}\NormalTok{()}
\end{Highlighting}
\end{Shaded}

\pandocbounded{\includegraphics[keepaspectratio]{Lab3_probability_files/figure-latex/unnamed-chunk-1-1.pdf}}

\subsection{Compared to What?}\label{compared-to-what}

We've shown that Kobe had some long shooting streaks, but are they long
enough to support the belief that he had a hot hand? What can we compare
them to?

To answer these questions, let's return to the idea of
\emph{independence}. Two processes are independent if the outcome of one
process doesn't effect the outcome of the second. If each shot that a
player takes is an independent process, having made or missed your first
shot will not affect the probability that you will make or miss your
second shot.

A shooter with a hot hand will have shots that are \emph{not}
independent of one another. Specifically, if the shooter makes his first
shot, the hot hand model says he will have a \emph{higher} probability
of making his second shot.

Let's suppose for a moment that the hot hand model is valid for Kobe.
During his career, the percentage of time Kobe makes a basket (i.e.~his
shooting percentage) is about 45\%, or in probability notation,

\[ P(\textrm{shot 1 = H}) = 0.45 \]

If he makes the first shot and has a hot hand (\emph{not} independent
shots), then the probability that he makes his second shot would go up
to, let's say, 60\%,

\[ P(\textrm{shot 2 = H} \, | \, \textrm{shot 1 = H}) = 0.60 \]

As a result of these increased probabilites, you'd expect Kobe to have
longer streaks. Compare this to the skeptical perspective where Kobe
does \emph{not} have a hot hand, where each shot is independent of the
next. If he hit his first shot, the probability that he makes the second
is still 0.45.

\[ P(\textrm{shot 2 = H} \, | \, \textrm{shot 1 = H}) = 0.45 \]

In other words, making the first shot did nothing to effect the
probability that he'd make his second shot. If Kobe's shots are
independent, then he'd have the same probability of hitting every shot
regardless of his past shots: 45\%.

Now that we've phrased the situation in terms of independent shots,
let's return to the question: how do we tell if Kobe's shooting streaks
are long enough to indicate that he has a hot hand? We can compare his
streak lengths to someone without a hot hand: an independent shooter.

\subsection{Simulations in R}\label{simulations-in-r}

While we don't have any data from a shooter we know to have independent
shots, that sort of data is very easy to simulate in R. In a simulation,
you set the ground rules of a random process and then the computer uses
random numbers to generate an outcome that adheres to those rules. As a
simple example, you can simulate flipping a fair coin with the
following.

\begin{Shaded}
\begin{Highlighting}[]
\NormalTok{coin\_outcomes }\OtherTok{\textless{}{-}} \FunctionTok{c}\NormalTok{(}\StringTok{"heads"}\NormalTok{, }\StringTok{"tails"}\NormalTok{)}
\FunctionTok{sample}\NormalTok{(coin\_outcomes, }\AttributeTok{size =} \DecValTok{1}\NormalTok{, }\AttributeTok{replace =} \ConstantTok{TRUE}\NormalTok{)}
\end{Highlighting}
\end{Shaded}

\begin{verbatim}
## [1] "heads"
\end{verbatim}

The vector \texttt{coin\_outcomes} can be thought of as a hat with two
slips of paper in it: one slip says \texttt{heads} and the other says
\texttt{tails}. The function \texttt{sample} draws one slip from the hat
and tells us if it was a head or a tail.

Run the second command listed above several times. Just like when
flipping a coin, sometimes you'll get a heads, sometimes you'll get a
tails, but in the long run, you'd expect to get roughly equal numbers of
each.

If you wanted to simulate flipping a fair coin 100 times, you could
either run the function 100 times or, more simply, adjust the
\texttt{size} argument, which governs how many samples to draw (the
\texttt{replace\ =\ TRUE} argument indicates we put the slip of paper
back in the hat before drawing again). Save the resulting vector of
heads and tails in a new object called \texttt{sim\_fair\_coin}.

\begin{Shaded}
\begin{Highlighting}[]
\NormalTok{sim\_fair\_coin }\OtherTok{\textless{}{-}} \FunctionTok{sample}\NormalTok{(coin\_outcomes, }\AttributeTok{size =} \DecValTok{100}\NormalTok{, }\AttributeTok{replace =} \ConstantTok{TRUE}\NormalTok{)}
\end{Highlighting}
\end{Shaded}

To view the results of this simulation, type the name of the object and
then use \texttt{table} to count up the number of heads and tails.

\begin{Shaded}
\begin{Highlighting}[]
\NormalTok{sim\_fair\_coin}
\end{Highlighting}
\end{Shaded}

\begin{verbatim}
##   [1] "tails" "tails" "tails" "heads" "heads" "tails" "tails" "tails" "tails"
##  [10] "heads" "tails" "heads" "tails" "heads" "tails" "heads" "heads" "heads"
##  [19] "tails" "heads" "tails" "tails" "heads" "heads" "heads" "tails" "tails"
##  [28] "tails" "tails" "tails" "tails" "tails" "heads" "heads" "heads" "heads"
##  [37] "heads" "tails" "tails" "heads" "heads" "tails" "tails" "tails" "heads"
##  [46] "tails" "tails" "tails" "heads" "tails" "heads" "heads" "heads" "tails"
##  [55] "tails" "tails" "tails" "tails" "heads" "heads" "tails" "heads" "heads"
##  [64] "heads" "heads" "heads" "tails" "heads" "heads" "heads" "tails" "tails"
##  [73] "tails" "heads" "tails" "tails" "tails" "tails" "tails" "heads" "heads"
##  [82] "heads" "tails" "heads" "tails" "tails" "heads" "heads" "heads" "tails"
##  [91] "tails" "heads" "heads" "heads" "tails" "tails" "heads" "tails" "heads"
## [100] "heads"
\end{verbatim}

\begin{Shaded}
\begin{Highlighting}[]
\FunctionTok{table}\NormalTok{(sim\_fair\_coin)}
\end{Highlighting}
\end{Shaded}

\begin{verbatim}
## sim_fair_coin
## heads tails 
##    48    52
\end{verbatim}

Since there are only two elements in \texttt{coin\_outcomes}, the
probability that we ``flip'' a coin and it lands heads is 0.5. Say we're
trying to simulate an unfair coin that we know only lands heads 20\% of
the time. We can adjust for this by adding an argument called
\texttt{prob}, which provides a vector of two probability weights.

\begin{Shaded}
\begin{Highlighting}[]
\NormalTok{sim\_unfair\_coin }\OtherTok{\textless{}{-}} \FunctionTok{sample}\NormalTok{(coin\_outcomes, }\AttributeTok{size =} \DecValTok{100}\NormalTok{, }\AttributeTok{replace =} \ConstantTok{TRUE}\NormalTok{, }
                          \AttributeTok{prob =} \FunctionTok{c}\NormalTok{(}\FloatTok{0.2}\NormalTok{, }\FloatTok{0.8}\NormalTok{))}
\end{Highlighting}
\end{Shaded}

\texttt{prob=c(0.2,\ 0.8)} indicates that for the two elements in the
\texttt{outcomes} vector, we want to select the first one,
\texttt{heads}, with probability 0.2 and the second one, \texttt{tails}
with probability 0.8. Another way of thinking about this is to think of
the outcome space as a bag of 10 chips, where 2 chips are labeled
``head'' and 8 chips ``tail''. Therefore at each draw, the probability
of drawing a chip that says ``head''\,'' is 20\%, and ``tail'' is 80\%.

\begin{enumerate}
\def\labelenumi{\arabic{enumi}.}
\setcounter{enumi}{2}
\tightlist
\item
  In your simulation of flipping the unfair coin 100 times, how many
  flips came up heads? Include the code for sampling the unfair coin in
  your response. Since the markdown file will run the code, and generate
  a new sample each time you \emph{Knit} it, you should also ``set a
  seed'' \textbf{before} you sample. Read more about setting a seed
  below.
\end{enumerate}

\textbf{Insert your answer here}by flippping the unfair coin 100 times
heads come up 14 times

\begin{Shaded}
\begin{Highlighting}[]
\FunctionTok{set.seed}\NormalTok{(}\DecValTok{1987}\NormalTok{)}
\NormalTok{sim\_unfair\_coin }\OtherTok{\textless{}{-}} \FunctionTok{sample}\NormalTok{(coin\_outcomes, }\AttributeTok{size =} \DecValTok{100}\NormalTok{, }\AttributeTok{replace =} \ConstantTok{TRUE}\NormalTok{, }\AttributeTok{prob =} \FunctionTok{c}\NormalTok{(}\FloatTok{0.2}\NormalTok{, }\FloatTok{0.8}\NormalTok{))}
\NormalTok{sim\_unfair\_coin}
\end{Highlighting}
\end{Shaded}

\begin{verbatim}
##   [1] "tails" "heads" "tails" "tails" "heads" "heads" "tails" "tails" "tails"
##  [10] "tails" "tails" "tails" "tails" "tails" "tails" "tails" "tails" "tails"
##  [19] "heads" "tails" "tails" "tails" "tails" "tails" "tails" "tails" "tails"
##  [28] "tails" "tails" "tails" "tails" "tails" "tails" "tails" "tails" "tails"
##  [37] "tails" "tails" "tails" "heads" "tails" "tails" "tails" "tails" "tails"
##  [46] "tails" "tails" "tails" "heads" "tails" "tails" "tails" "tails" "tails"
##  [55] "tails" "tails" "tails" "tails" "tails" "tails" "tails" "tails" "tails"
##  [64] "tails" "tails" "tails" "tails" "tails" "heads" "tails" "tails" "tails"
##  [73] "tails" "heads" "tails" "tails" "tails" "tails" "tails" "heads" "tails"
##  [82] "tails" "tails" "heads" "tails" "tails" "heads" "tails" "tails" "heads"
##  [91] "heads" "tails" "tails" "heads" "tails" "tails" "tails" "tails" "tails"
## [100] "tails"
\end{verbatim}

\begin{Shaded}
\begin{Highlighting}[]
\FunctionTok{table}\NormalTok{(sim\_unfair\_coin)}
\end{Highlighting}
\end{Shaded}

\begin{verbatim}
## sim_unfair_coin
## heads tails 
##    14    86
\end{verbatim}

\phantomsection\label{boxedtext}
\textbf{A note on setting a seed:} Setting a seed will cause R to select
the same sample each time you knit your document. This will make sure
your results don't change each time you knit, and it will also ensure
reproducibility of your work (by setting the same seed it will be
possible to reproduce your results). You can set a seed like this:

\begin{Shaded}
\begin{Highlighting}[]
\FunctionTok{set.seed}\NormalTok{(}\DecValTok{35797}\NormalTok{)                 }\CommentTok{\# make sure to change the seed}
\end{Highlighting}
\end{Shaded}

The number above is completely arbitraty. If you need inspiration, you
can use your ID, birthday, or just a random string of numbers. The
important thing is that you use each seed only once in a document.
Remember to do this \textbf{before} you sample in the exercise above.

In a sense, we've shrunken the size of the slip of paper that says
``heads'', making it less likely to be drawn, and we've increased the
size of the slip of paper saying ``tails'', making it more likely to be
drawn. When you simulated the fair coin, both slips of paper were the
same size. This happens by default if you don't provide a \texttt{prob}
argument; all elements in the \texttt{outcomes} vector have an equal
probability of being drawn.

If you want to learn more about \texttt{sample} or any other function,
recall that you can always check out its help file.

\begin{Shaded}
\begin{Highlighting}[]
\NormalTok{?sample}
\end{Highlighting}
\end{Shaded}

\subsection{Simulating the Independent
Shooter}\label{simulating-the-independent-shooter}

Simulating a basketball player who has independent shots uses the same
mechanism that you used to simulate a coin flip. To simulate a single
shot from an independent shooter with a shooting percentage of 50\% you
can type

\begin{Shaded}
\begin{Highlighting}[]
\NormalTok{shot\_outcomes }\OtherTok{\textless{}{-}} \FunctionTok{c}\NormalTok{(}\StringTok{"H"}\NormalTok{, }\StringTok{"M"}\NormalTok{)}
\NormalTok{sim\_basket }\OtherTok{\textless{}{-}} \FunctionTok{sample}\NormalTok{(shot\_outcomes, }\AttributeTok{size =} \DecValTok{1}\NormalTok{, }\AttributeTok{replace =} \ConstantTok{TRUE}\NormalTok{, }\AttributeTok{prob =} \FunctionTok{c}\NormalTok{(}\FloatTok{0.5}\NormalTok{, }\FloatTok{0.5}\NormalTok{))}
\NormalTok{sim\_basket}
\end{Highlighting}
\end{Shaded}

\begin{verbatim}
## [1] "H"
\end{verbatim}

\begin{Shaded}
\begin{Highlighting}[]
\FunctionTok{table}\NormalTok{(sim\_basket)}
\end{Highlighting}
\end{Shaded}

\begin{verbatim}
## sim_basket
## H 
## 1
\end{verbatim}

To make a valid comparison between Kobe and your simulated independent
shooter, you need to align both their shooting percentage and the number
of attempted shots.

\begin{enumerate}
\def\labelenumi{\arabic{enumi}.}
\setcounter{enumi}{3}
\tightlist
\item
  What change needs to be made to the \texttt{sample} function so that
  it reflects a shooting percentage of 45\%? Make this adjustment, then
  run a simulation to sample 133 shots. Assign the output of this
  simulation to a new object called \texttt{sim\_basket}.
\end{enumerate}

\textbf{Insert your answer here} the changes that need to be made is to
add a prob \#we get hits =47 and miss = 86 randomly output

\begin{Shaded}
\begin{Highlighting}[]
\NormalTok{shot\_outcomes }\OtherTok{\textless{}{-}} \FunctionTok{c}\NormalTok{(}\StringTok{"H"}\NormalTok{, }\StringTok{"M"}\NormalTok{)}
\NormalTok{sim\_basket }\OtherTok{\textless{}{-}} \FunctionTok{sample}\NormalTok{(shot\_outcomes, }\AttributeTok{size =}  \DecValTok{133}\NormalTok{, }\AttributeTok{replace =} \ConstantTok{TRUE}\NormalTok{, }\AttributeTok{prob =} \FunctionTok{c}\NormalTok{(}\FloatTok{0.45}\NormalTok{, }\FloatTok{0.55}\NormalTok{))}
\NormalTok{sim\_basket}
\end{Highlighting}
\end{Shaded}

\begin{verbatim}
##   [1] "M" "H" "M" "H" "H" "M" "H" "H" "M" "M" "H" "H" "H" "H" "H" "H" "M" "M"
##  [19] "M" "M" "H" "H" "M" "M" "H" "M" "M" "M" "H" "H" "H" "M" "H" "M" "H" "M"
##  [37] "M" "M" "M" "M" "M" "H" "M" "H" "H" "H" "M" "M" "H" "H" "M" "H" "H" "H"
##  [55] "H" "M" "M" "M" "M" "M" "H" "M" "H" "H" "H" "H" "H" "M" "M" "M" "H" "H"
##  [73] "M" "M" "M" "H" "M" "H" "H" "M" "H" "M" "M" "H" "M" "H" "M" "M" "H" "M"
##  [91] "H" "H" "M" "M" "M" "H" "H" "M" "M" "H" "H" "H" "M" "M" "M" "H" "H" "H"
## [109] "M" "M" "H" "M" "H" "M" "M" "M" "M" "M" "M" "M" "M" "H" "H" "M" "M" "M"
## [127] "H" "H" "H" "H" "M" "M" "M"
\end{verbatim}

\begin{Shaded}
\begin{Highlighting}[]
\FunctionTok{table}\NormalTok{(sim\_basket)}
\end{Highlighting}
\end{Shaded}

\begin{verbatim}
## sim_basket
##  H  M 
## 62 71
\end{verbatim}

Note that we've named the new vector \texttt{sim\_basket}, the same name
that we gave to the previous vector reflecting a shooting percentage of
50\%. In this situation, R overwrites the old object with the new one,
so always make sure that you don't need the information in an old vector
before reassigning its name.

With the results of the simulation saved as \texttt{sim\_basket}, you
have the data necessary to compare Kobe to our independent shooter.

Both data sets represent the results of 133 shot attempts, each with the
same shooting percentage of 45\%. We know that our simulated data is
from a shooter that has independent shots. That is, we know the
simulated shooter does not have a hot hand.

\begin{center}\rule{0.5\linewidth}{0.5pt}\end{center}

\subsection{More Practice}\label{more-practice}

\subsubsection{Comparing Kobe Bryant to the Independent
Shooter}\label{comparing-kobe-bryant-to-the-independent-shooter}

\begin{enumerate}
\def\labelenumi{\arabic{enumi}.}
\setcounter{enumi}{4}
\tightlist
\item
  Using \texttt{calc\_streak}, compute the streak lengths of
  \texttt{sim\_basket}, and save the results in a data frame called
  \texttt{sim\_streak}.
\end{enumerate}

\textbf{Insert your answer here} computing the streak length

\begin{Shaded}
\begin{Highlighting}[]
\NormalTok{sim\_streak }\OtherTok{\textless{}{-}}\NormalTok{ openintro}\SpecialCharTok{::}\FunctionTok{calc\_streak}\NormalTok{(sim\_basket)}
\NormalTok{sim\_streak}
\end{Highlighting}
\end{Shaded}

\begin{verbatim}
##    length
## 1       0
## 2       1
## 3       2
## 4       2
## 5       0
## 6       6
## 7       0
## 8       0
## 9       0
## 10      2
## 11      0
## 12      1
## 13      0
## 14      0
## 15      3
## 16      1
## 17      1
## 18      0
## 19      0
## 20      0
## 21      0
## 22      0
## 23      1
## 24      3
## 25      0
## 26      2
## 27      4
## 28      0
## 29      0
## 30      0
## 31      0
## 32      1
## 33      5
## 34      0
## 35      0
## 36      2
## 37      0
## 38      0
## 39      1
## 40      2
## 41      1
## 42      0
## 43      1
## 44      1
## 45      0
## 46      1
## 47      2
## 48      0
## 49      0
## 50      2
## 51      0
## 52      3
## 53      0
## 54      0
## 55      3
## 56      0
## 57      1
## 58      1
## 59      0
## 60      0
## 61      0
## 62      0
## 63      0
## 64      0
## 65      0
## 66      2
## 67      0
## 68      0
## 69      4
## 70      0
## 71      0
## 72      0
\end{verbatim}

\begin{enumerate}
\def\labelenumi{\arabic{enumi}.}
\setcounter{enumi}{5}
\tightlist
\item
  Describe the distribution of streak lengths. What is the typical
  streak length for this simulated independent shooter with a 45\%
  shooting percentage? How long is the player's longest streak of
  baskets in 133 shots? Make sure to include a plot in your answer.
\end{enumerate}

\textbf{Insert your answer here} The distribution of simulated's streak
lengths is right-skewed. Most of the streaks are short (0 or 1), and the
frequency decreases as the streak length increases. This means short
streaks occurred much more often than long streaks. the typical streak
length is 1 the player longest streak is 9

\begin{Shaded}
\begin{Highlighting}[]
\FunctionTok{ggplot}\NormalTok{(}\AttributeTok{data =}\NormalTok{ sim\_streak, }\FunctionTok{aes}\NormalTok{( }\AttributeTok{x =}\NormalTok{ length)) }\SpecialCharTok{+} \FunctionTok{geom\_bar}\NormalTok{()}
\end{Highlighting}
\end{Shaded}

\pandocbounded{\includegraphics[keepaspectratio]{Lab3_probability_files/figure-latex/unnamed-chunk-5-1.pdf}}

\begin{enumerate}
\def\labelenumi{\arabic{enumi}.}
\setcounter{enumi}{6}
\tightlist
\item
  If you were to run the simulation of the independent shooter a second
  time, how would you expect its streak distribution to compare to the
  distribution from the question above? Exactly the same? Somewhat
  similar? Totally different? Explain your reasoning.
\end{enumerate}

\textbf{Insert your answer here} If I were to run the simulation a
second time, I would expect the streak distribution to be somewhat
similar but not exactly the same. Because the simulation uses random
sampling, each run produces slightly different results. However, since
the shooting percentage remains 45\%, the overall shape of the
distribution should remain right-skewed with mostly short streaks.

\begin{enumerate}
\def\labelenumi{\arabic{enumi}.}
\setcounter{enumi}{7}
\tightlist
\item
  How does Kobe Bryant's distribution of streak lengths compare to the
  distribution of streak lengths for the simulated shooter? Using this
  comparison, do you have evidence that the hot hand model fits Kobe's
  shooting patterns? Explain.
\end{enumerate}

\textbf{Insert your answer here} Kobe Bryant's distribution of streak
lengths is very similar to that of the simulated independent shooter.
Both distributions are right-skewed, with most streaks being short and
only a few long streaks.

Since Kobe's longest streak and overall pattern do not appear
substantially different from what we would expect under independence,
there is no strong evidence supporting the hot hand model. His shooting
performance is consistent with an independent process. * * *

\end{document}
